%%% FIELD proof-sharp-three-facet-endpoints
Fix \(L\geq 2\), a compatible arm-marginal vector \(p\), and \(d\in\mathcal D\).  Write \(p^+=(p_{rj})_{r\in\mathcal D,\,1\leq j<L}\), the marginal coordinates retained in \(v_L\).  We first verify that the canonical rows give the complete inequality description needed below.

By \(\text{thm:dimension-and-staircase-facets}\), \(Q_{3,L}\) is full-dimensional in \(\mathbb R^{3L}\).  A facet therefore contains \(3L\) affinely independent vertices.  Those vertices have integral coordinates, so the one-dimensional nullspace of their \(3L\)-by-\((3L+1)\) augmented integer matrix has a nonzero rational generator.  Multiplying by a common denominator, dividing by the positive greatest common divisor, and choosing the valid orientation produces a unique primitive integer row \(h\) of \(\text{def:canonical-facet-set}\).  Conversely, the validity and tight-set rank conditions in that definition say exactly that every \(h\in\mathcal H_L\) supports a facet.  Hence \(\mathcal H_L\) is precisely the finite set of oriented primitive facet rows.

These facet rows describe all of \(Q_{3,L}\).  To see this directly, choose an interior point \(x_0\in Q_{3,L}\).  If \(x\notin Q_{3,L}\), compactness of the finite convex hull implies that the ray segment from \(x_0\) toward \(x\) has a first boundary point.  A boundary point of a full-dimensional polytope belongs to a facet.  The oriented slack of that facet is positive at \(x_0\), zero at the first boundary point, and negative farther along the same ray, in particular at \(x\).  Thus every exterior point violates a row in \(\mathcal H_L\).

For \(h_i\in\mathcal H_L\), write its inequality at coordinates \((\vartheta,p^+)\) as
\[
 a_i(p)+\beta_i^{\mathsf T}\vartheta\geq0,
 \qquad
 a_i(p)=\kappa_i+\sum_{r\in\mathcal D}\sum_{j=1}^{L-1}\alpha_{irj}p_{rj}.
\]
It follows that the marginal section is
\[
 K(p):=\{\vartheta\in\mathbb R^3:(\vartheta,p^+)\in Q_{3,L}\}
 =\{\vartheta:a_i(p)+\beta_i^{\mathsf T}\vartheta\geq0
       \text{ for every }i\in\mathcal H_L\}.
 \tag{1}
\]
Convex weights on the vertices \(v_L(u)\), \(u\in\mathcal U_L\), are exactly probability masses \(q\in\Delta(\mathcal U_L)\).  Fixing the marginal coordinates gives precisely the equalities in \(\text{def:coupling-fiber}\), including the level-zero equalities because each arm marginal sums to one, while the first three coordinates are \(\theta(q)\) by \(\text{def:unique-peak-indicators}\).  Therefore
\[
 K(p)=\{\theta(q):q\in\Gamma_L(p)\}.                 \tag{2}
\]
Compatibility makes this set nonempty, and finiteness of \(\mathcal U_L\) makes it compact.  Equations (1)--(2), consistently with the support-constrained primal representation in \(\text{thm:direct-support-transport-dual}\), give
\[
 \underline\theta_d(p)=\min_{\vartheta\in K(p)}e_d^{\mathsf T}\vartheta,
 \qquad
 \overline\theta_d(p)=\max_{\vartheta\in K(p)}e_d^{\mathsf T}\vartheta.       \tag{3}
\]

Apply \(\text{lem:finite-lp-strong-duality}\) to the first program in (3), written with the finitely many constraints in (1).  Its dual is
\[
 \max_{\lambda\geq0}
 \left\{-\sum_i\lambda_i a_i(p):
             \sum_i\lambda_i\beta_i=e_d\right\}.       \tag{4}
\]
The primal optimum exists by compactness, so the cited finite-dimensional duality result yields a dual optimizer and equality of the two values.  Weak duality is also visible equation by equation: for every feasible \(\vartheta\) and \(\lambda\),
\[
 0\leq\sum_i\lambda_i\{a_i(p)+\beta_i^{\mathsf T}\vartheta\}
 =\sum_i\lambda_i a_i(p)+e_d^{\mathsf T}\vartheta.     \tag{5}
\]
Let \(\vartheta^-\) and \(\lambda\) be primal and dual optimizers.  Equality of their objective values makes (5) an equality at \(\vartheta^-\).  Since every summand on its left is nonnegative,
\[
 \lambda_i>0\quad\Longrightarrow\quad
 a_i(p)+\beta_i^{\mathsf T}\vartheta^-=0.              \tag{6}
\]

We now sparsify this optimal certificate without changing its value.  If the normals on \(S=\operatorname{supp}(\lambda)\) are linearly dependent, choose a nonzero vector \(r=(r_i)_{i\in S}\) satisfying \(\sum_{i\in S}r_i\beta_i=0\), and reverse its sign if necessary so that some \(r_i>0\).  Set
\[
 t=\min_{i\in S:r_i>0}\frac{\lambda_i}{r_i}>0,
 \qquad
 \lambda'_i=\lambda_i-tr_i\quad(i\in S),
\]
with \(\lambda'_i=0\) off \(S\).  Then \(\lambda'\geq0\), at least one positive coefficient is removed, and \(\sum_i\lambda'_i\beta_i=e_d\).  By (6),
\[
 \sum_{i\in S}r_i a_i(p)
 =-\left(\sum_{i\in S}r_i\beta_i\right)^{\mathsf T}\vartheta^-=0,             \tag{7}
\]
so the objective in (4) is unchanged.  Iterating this finite support-reduction step leaves an optimal certificate with linearly independent normals.  There can be at most three such normals in \(\mathbb R^3\); the resulting certificate lies in \(\mathcal D_d^-\).  Conversely, (5) proves that every member of \(\mathcal D_d^-\) supplies a valid lower bound.  Thus
\[
 \underline\theta_d(p)=
 \max_{\lambda\in\mathcal D_d^-}-\sum_i\lambda_i a_i(p).                    \tag{8}
\]

For the upper endpoint, minimize \(-e_d^{\mathsf T}\vartheta\) over \(K(p)\).  The same argument gives a nonnegative dual optimizer \(\mu\) satisfying \(\sum_i\mu_i\beta_i=-e_d\), and its value identity is
\[
 \overline\theta_d(p)=
 \min_{\mu\in\mathcal D_d^+}\sum_i\mu_i a_i(p).                            \tag{9}
\]
The identical active-facet reduction proves that an optimizer in (9) has at most three linearly independent full-facet normals.  There are only finitely many supports of size at most three, and on each linearly independent support the equation \(\sum_i\rho_i\beta_i=\pm e_d\) has at most one solution.  Hence \(\mathcal D_d^-\) and \(\mathcal D_d^+\) are finite and computable from \(\mathcal H_L\).  Equations (2), (8), and (9) show equality with extrema over compatible latent laws, not merely outer bounds.  Together with \(\text{thm:randomization-identifies-marginals}\), they distinguish the causal target \(\theta(q)\) from the observed-law functionals in (8)--(9) and prove that those functionals are its sharp identified endpoints.

It remains to prove that the number three cannot be replaced uniformly by two.  Let \(L=2\), put \(x_r=p_{r1}\), and list the seven legal triples as
\[
 000,001,010,011,100,110,111;
\]
the excluded triple is \(101\) by \(\text{def:non-valley-support}\).  Solving the affine vertex equations gives
\[
\begin{aligned}
 q_{000}&=1-x_1-\theta_0-\theta_2,&
 q_{001}&=\theta_2,&
 q_{010}&=\theta_1,\\
 q_{011}&=x_1-x_0-\theta_1+\theta_0,&
 q_{100}&=\theta_0,\\
 q_{110}&=x_1-x_2-\theta_1+\theta_2,&
 q_{111}&=x_0+x_2-x_1-\theta_0-\theta_2+\theta_1.
\end{aligned}                                                                  \tag{10}
\]
Thus the affine vertex map is invertible.  Since \(Q_{3,2}\) has dimension six, it is the simplex generated by these seven vertices, and its seven canonical full facets are exactly the nonnegativity inequalities in (10).  Their \(\beta\)-normals, in the displayed order, are
\[
 (-1,0,-1),\ (0,0,1),\ (0,1,0),\ (1,-1,0),\ (1,0,0),\ (0,-1,1),\ (-1,1,-1). \tag{11}
\]
Only \((1,-1,0)\) and \((0,-1,1)\) have a negative second coordinate.  If a nonnegative combination with support at most two used only the first of these, cancellation of its positive first coordinate would require either \((-1,0,-1)\) or \((-1,1,-1)\), leaving an uncancelled negative third coordinate.  The symmetric argument applies if it used only the second.  If it used both, the sum is \((r,-r-s,s)\) with \(r,s>0\), whose first and third coordinates cannot both vanish.  Therefore no nonnegative combination of at most two normals in (11) equals \(-e_1\).  In contrast,
\[
 \frac12\bigl\{(-1,0,-1)+(1,-1,0)+(0,-1,1)\bigr\}=-e_1.                    \tag{12}
\]
Combining the corresponding three affine inequalities gives
\[
 \theta_1\leq\frac{1+x_1-x_0-x_2}{2}.                                      \tag{13}
\]
At \(x_0=x_1=x_2=1/2\), (13) gives \(\theta_1\leq1/4\).  The legal coupling assigning mass \(1/4\) to each of \(001,010,100,111\) has all three arm marginals equal to \(1/2\) and has \(\theta_1=1/4\).  Hence (13) is sharp there, while (11)--(12) prove that every full-facet certificate for this upper endpoint needs at least three normals.  This completes both the sharp endpoint formulas and the uniform necessity claim.

%%% FIELD proof-endpoint-attaining-mixtures
Fix \(L\geq2\), a compatible \(p\), a dose \(d\in\mathcal D\), and either endpoint objective.  Index columns by \(u\in\mathcal U_L\).  Let \(M\) be the \((3L-2)\)-by-\(|\mathcal U_L|\) matrix with a row \(\varnothing\) and rows \((r,j)\in\mathcal D\times\{1,\ldots,L-1\}\), defined by
\[
 M_{\varnothing,u}=1,
 \qquad
 M_{(r,j),u}=\mathbf 1\{u_r=j\},
\]
and set
\[
 b(p)=\bigl(1,(p_{rj})_{r\in\mathcal D,\,1\leq j<L}\bigr)^{\mathsf T}.
\]
Then \(\text{def:coupling-fiber}\) is equivalently
\[
 \Gamma_L(p)=\{q\in\mathbb R_+^{\mathcal U_L}:Mq=b(p)\}.                    \tag{14}
\]
Indeed, a member of the fiber satisfies these equations.  Conversely, the mass row in (14) gives \(\sum_uq(u)=1\), and for every arm \(r\), the omitted level-zero equation follows from
\[
 \sum_{u:u_r=0}q(u)
 =1-\sum_{j=1}^{L-1}\sum_{u:u_r=j}q(u)
 =1-\sum_{j=1}^{L-1}p_{rj}=p_{r0},                                         \tag{15}
\]
where the last equality uses that the compatible arm marginal belongs to \(\Delta^{L-1}\).

The feasible set in (14) is nonempty by compatibility, closed, and contained in the finite probability simplex, hence compact.  The map
\[
 q\longmapsto\theta_d(q)=\sum_{u\in\mathcal U_L}B_d(u)q(u)
\]
is linear and continuous, so both endpoint extrema are attained.  For either objective choose an optimizer \(q^\star\) with minimum support and put \(S=\operatorname{supp}(q^\star)\).  If the columns \((M_u)_{u\in S}\) were linearly dependent, there would be a nonzero vector \(z\), supported on \(S\), with \(Mz=0\).  For sufficiently small \(\varepsilon>0\), both \(q^\star+\varepsilon z\) and \(q^\star-\varepsilon z\) would be nonnegative and feasible.  Optimality of both signed perturbations forces
\[
 \sum_{u\in\mathcal U_L}B_d(u)z(u)=0.                                      \tag{16}
\]
The mass row of \(Mz=0\) gives \(\sum_uz(u)=0\), so a nonzero \(z\) has both signs.  Therefore
\[
 t=\min_{u\in S:z(u)<0}\frac{q^\star(u)}{-z(u)}>0
\]
makes \(q^\star+tz\) feasible, optimal by (16), and supported on a strict subset of \(S\), a contradiction.  The columns on \(S\) are therefore linearly independent, and
\[
 |\operatorname{supp}(q^\star)|\leq\operatorname{rank}(M)\leq3L-2.          \tag{17}
\]

We next make the promised construction from an active facet certificate explicit.  For \(h_i\in\mathcal H_L\), write
\[
 a_i(p)=\kappa_i+\sum_{r\in\mathcal D}\sum_{j=1}^{L-1}\alpha_{irj}p_{rj}.
\]
For every \(q\in\Gamma_L(p)\), the marginal equalities, the definition of \(\theta(q)\), and the vertex slack in \(\text{def:canonical-facet-set}\) give
\[
 a_i(p)+\beta_i^{\mathsf T}\theta(q)
 =\sum_{u\in\mathcal U_L}s_{h_i}(u)q(u)\geq0.                              \tag{18}
\]
For a lower endpoint, \(\text{thm:sharp-three-facet-endpoints}\) supplies an optimal \(\lambda\in\mathcal D_d^-\) with
\[
 \sum_i\lambda_i\beta_i=e_d,
 \qquad
 \underline\theta_d(p)=-\sum_i\lambda_i a_i(p),
 \qquad \lambda_i\geq0.                                                     \tag{19}
\]
For any lower-endpoint optimizer, summing (18) with these weights yields
\[
 0=\underline\theta_d(p)+\sum_i\lambda_i a_i(p)
 =\sum_i\lambda_i\sum_{u\in\mathcal U_L}s_{h_i}(u)q^\star(u).             \tag{20}
\]
Every term is nonnegative.  Hence \(\lambda_i>0\) and \(q^\star(u)>0\) imply \(s_{h_i}(u)=0\).  For an upper endpoint the same theorem supplies an optimal \(\mu\in\mathcal D_d^+\) satisfying
\[
 \sum_i\mu_i\beta_i=-e_d,
 \qquad
 \overline\theta_d(p)=\sum_i\mu_i a_i(p),
 \qquad \mu_i\geq0,                                                        \tag{21}
\]
and summing (18) gives the same conclusion from
\[
 0=\sum_i\mu_i a_i(p)-\overline\theta_d(p)
 =\sum_i\mu_i\sum_{u\in\mathcal U_L}s_{h_i}(u)q^\star(u).                \tag{22}
\]

Let \(\rho\) denote the relevant certificate, either \(\lambda\) or \(\mu\), and define
\[
 T_\rho=\{u\in\mathcal U_L:s_{h_i}(u)=0
             \text{ for every }i\text{ with }\rho_i>0\}.                  \tag{23}
\]
Equations (20)--(22) show that
\[
 M_{T_\rho}q_{T_\rho}=b(p),\qquad q_{T_\rho}\geq0                         \tag{24}
\]
is feasible.  Conversely, extend any solution of (24) by zero outside \(T_\rho\).  Equation (14) puts it in \(\Gamma_L(p)\), while (18) summed with the certificate weights gives equality in (19) or (21).  It therefore attains the specified endpoint.

This yields a finite constructive procedure.  Enumerate subsets \(S\subseteq T_\rho\) with \(|S|\leq3L-2\), solve
\[
 M_Sq_S=b(p),\qquad q_S\geq0,                                               \tag{25}
\]
and accept a solution whose columns are linearly independent.  Feasibility of (24) and the support-elimination argument leading to (17), applied within \(T_\rho\), guarantee that the enumeration succeeds.  If \(p\) is rational, then \(b(p)\) is rational and \(M\) is integral.  On the linearly independent support of the active certificate, the equation \(\sum_i\rho_i\beta_i=\pm e_d\) has a unique solution obtained from a nonsingular integer minor, so its weights are rational; the tight set (23) is exactly decidable because every slack is integral.  For an accepted \(S\), choose \(|S|\) rows \(R\) such that the square integer matrix \(M_{R,S}\) is nonsingular.  Then
\[
 q_S=M_{R,S}^{-1}b_R(p)                                                     \tag{26}
\]
has rational entries.  Exact rational elimination checks the remaining equations and nonnegativity in (25).  Thus complementary slackness with an active certificate constructs a rational endpoint-attaining mixture whenever \(p\) is rational.

Finally, (17) is uniformly sharp.  Take \(L=2\) and \(p_{0,1}=p_{1,1}=p_{2,1}=1/2\).  The only deleted binary triple is \(101\).  Write \(x_{abc}=q(a,b,c)\) for the seven legal triples.  The marginal equations are
\[
\begin{aligned}
 x_{100}+x_{110}+x_{111}&=\frac12,\\
 x_{010}+x_{011}+x_{110}+x_{111}&=\frac12,\\
 x_{001}+x_{011}+x_{111}&=\frac12.
\end{aligned}                                                               \tag{27}
\]
Adding the first and third equations and subtracting total mass one gives
\[
 x_{000}=x_{111}-x_{010}.                                                    \tag{28}
\]
Nonnegativity implies \(x_{111}\geq x_{010}\).  By \(\text{def:unique-peak-indicators}\), \(B_1(u)=1\) only at \(u=010\), so \(\theta_1(q)=x_{010}\).  The middle equation in (27) and (28) imply
\[
 \frac12=x_{010}+x_{011}+x_{110}+x_{111}\geq2x_{010},
\]
and hence \(\overline\theta_1(p)\leq1/4\).  The coupling
\[
 x_{001}=x_{010}=x_{100}=x_{111}=\frac14,
 \qquad x_{000}=x_{011}=x_{110}=0                                          \tag{29}
\]
is legal, satisfies (27), and attains \(\theta_1=1/4\).  If equality holds, then (27)--(28) force \(x_{111}=x_{010}=1/4\) and \(x_{011}=x_{110}=0\); the first and third equations force \(x_{100}=x_{001}=1/4\), and (28) forces \(x_{000}=0\).  Thus (29) is the unique optimizer and every optimizer has support \(4=3L-2\).  Consequently no smaller uniform support bound is possible.

%%% FIELD proof-fixed-l-endpoint-lipschitz-rate
Fix \(L\geq2\).  By \(\text{thm:sharp-three-facet-endpoints}\), each certificate set \(\mathcal D_d^-\) and \(\mathcal D_d^+\) is finite and nonempty, and its affine maximum or minimum agrees with the corresponding sharp endpoint on every compatible marginal vector.

For a lower certificate \(\lambda\), the gradient of its affine contrast with respect to \(x=(x_{kj})_{k\in\mathcal D,\,1\leq j<L}\) is
\[
 -\left(\sum_i\lambda_i\alpha_{ikj}\right)_{k,j},
\]
and for an upper certificate \(\mu\) it is the same array without the minus sign.  By the definition of the dual norm, every such affine contrast \(f\) obeys
\[
 |f(x)-f(x')|
 =|\langle\nabla f,x-x'\rangle|
 \leq\|\nabla f\|_{r^*}\|x-x'\|_r.                                      \tag{30}
\]
For any two finite real families indexed by the same set,
\[
 \left|\max_t a_t-\max_t b_t\right|\leq\max_t|a_t-b_t|,
 \qquad
 \left|\min_t a_t-\min_t b_t\right|\leq\max_t|a_t-b_t|.                 \tag{31}
\]
Indeed, choosing an index maximizing the first family bounds the first directed difference, and exchanging the two families bounds the reverse difference; minima follow after multiplying by \(-1\).  Applying (30)--(31) to the certificate lists gives
\[
 \max_{d\in\mathcal D,\,s\in\{-,+\}}
 |F_{d,L}^s(x)-F_{d,L}^s(x')|
 \leq C_{L,r}\|x-x'\|_r.                                                  \tag{32}
\]
The constant is finite because the lists are finite.  It is exactly computable from \(\mathcal H_L\): enumerate supports of at most three linearly independent normals, solve the corresponding systems \(\sum_i\rho_i\beta_i=\pm e_d\), retain the nonnegative solutions, and maximize the displayed dual norms.  The endpoint equalities on compatible \(p\) are precisely those already proved in \(\text{thm:sharp-three-facet-endpoints}\).

We next prove the finite-sample probability bound.  Put \(\pi_{\min}=\min_{d\in\mathcal D}\pi_d>0\), using \(\text{ass:positive-randomization}\), and define
\[
 E_N=\bigcap_{d\in\mathcal D}\{N_d\geq N\pi_d/2\}.
\]
Under \(\text{ass:iid-sampling}\), \(N_d\) is binomial with mean \(N\pi_d\) and variance \(N\pi_d(1-\pi_d)\).  Since
\[
 \mathbf 1\{|N_d-N\pi_d|\geq N\pi_d/2\}
 \leq \frac{4(N_d-N\pi_d)^2}{N^2\pi_d^2},
\]
taking expectations and then using \(\mathbf 1_{\cup_d E_d}\leq\sum_d\mathbf 1_{E_d}\) gives
\[
 P(E_N^c)
 \leq\sum_{d\in\mathcal D}\frac{4(1-\pi_d)}{N\pi_d}
 \leq\frac{12}{N\pi_{\min}}.                                              \tag{33}
\]

Conditional on the arm counts, the outcome counts in arm \(d\) are multinomial with cell probabilities \((p_{dj})_{j\in\mathcal Y_L}\).  On \(E_N\), all denominators are positive and
\[
\begin{aligned}
 \mathbb E\!\left(\|\widehat p-p\|_2^2\mid(N_d)_{d\in\mathcal D}\right)
 &=\sum_{d\in\mathcal D}\frac1{N_d}
   \sum_{j=1}^{L-1}p_{dj}(1-p_{dj})\\
 &\leq\sum_{d\in\mathcal D}\frac1{N_d}
 \leq\frac6{N\pi_{\min}}.                                                 \tag{34}
\end{aligned}
\]
Only coordinate variances enter the squared Euclidean norm; multinomial covariances therefore do not enter (34).  Pointwise,
\[
 \mathbf 1\!\left\{\|\widehat p-p\|_2>
       \frac{M}{\sqrt{N\pi_{\min}}}\right\}
 \leq\frac{N\pi_{\min}}{M^2}\|\widehat p-p\|_2^2.
\]
Taking conditional expectations on count vectors in \(E_N\), using (34), and averaging over the counts yields
\[
 P\!\left\{\|\widehat p-p\|_2>
       \frac{M}{\sqrt{N\pi_{\min}}},\ E_N\right\}
 \leq\frac6{M^2}.                                                          \tag{35}
\]
On \(E_N\) no arm is empty, so the arbitrary simplex convention for an empty arm is irrelevant.  Combining (32) with \(r=2\), (33), and (35) proves
\[
 P\!\left\{
 \max_{d,s}|F_{d,L}^s(\widehat p)-F_{d,L}^s(p)|>
 \frac{C_{L,2}M}{\sqrt{N\pi_{\min}}}
 \right\}
 \leq\frac6{M^2}+\frac{12}{N\pi_{\min}}.                                 \tag{36}
\]
If \(N\pi_{\min}\to\infty\), (36) is exactly the asserted \(O_P((N\pi_{\min})^{-1/2})\) statement.  At a fixed law, positivity fixes \(\pi_{\min}>0\), giving the \(N^{-1/2}\) rate and the stated constant dependence.  This proof uses the endpoint's finite affine representation and elementary multinomial second moments; it invokes no rate transfer from a different transport objective.

%%% FIELD proof-finite-sample-support-honesty
Fix \(L\geq2\), \(N\geq1\), \(\alpha\in(0,1)\), and \(P\in\mathcal P_L\).  Split the observations as in \(\text{def:inference-handle}\), with pilot size \(n_0=\lfloor N^{2/3}\rfloor\) and main size \(n_1=N-n_0\).  The subsamples are independent by \(\text{ass:iid-sampling}\).  Conditional on the arm counts in either subsample, random assignment, consistency, and \(\text{thm:randomization-identifies-marginals}\) imply that the three within-arm outcome-count vectors are independent multinomials with respective cell probabilities \((p_{dj})_{j\in\mathcal Y_L}\).

Set \(\delta_N=\min\{\alpha/2,N^{-2}\}\) and \(\gamma_N=\delta_N/(3L)\).  For \(X\sim\operatorname{Bin}(n,\rho)\) and an observed count \(k\), define
\[
 C_{n,k}(\gamma_N)=\left\{\rho\in[0,1]:
 P_\rho(X\leq k)\geq\frac{\gamma_N}{2},\quad
 P_\rho(X\geq k)\geq\frac{\gamma_N}{2}\right\},                         \tag{37}
\]
and set \(C_{0,0}(\gamma_N)=[0,1]\).  For \(n>0\), binomial tail probabilities are polynomials in \(\rho\), so (37) is closed.  At the true \(\rho\), the event \(P_\rho(X\leq X_{\mathrm{obs}})<\gamma_N/2\) is empty or a lower set \(\{X_{\mathrm{obs}}\leq k_0\}\), and its probability is then the corresponding lower tail, less than \(\gamma_N/2\).  The analogous upper-tail event has probability less than \(\gamma_N/2\).  Since the indicator of a union is bounded by the sum of its two indicators,
\[
 P_\rho\{\rho\in C_{n,X_{\mathrm{obs}}}(\gamma_N)\}\geq1-\gamma_N.        \tag{38}
\]
This argument includes \(\rho=0\) and \(\rho=1\).

For pilot arm \(d\), intersect its \(L\) cell sets with \(\Delta^{L-1}\).  If the arm is empty or the intersection is empty, replace it by the full simplex.  Let \(R_N\) be the Cartesian product of the three resulting nonempty compact sets.  Conditional on each pilot arm count, every cell count is marginally binomial.  Applying (38) to all \(3L\) cells and using the pointwise inequality between the indicator of a union and the sum of its indicators, then averaging over the arm counts, gives
\[
 P_p\{p\in R_N\}\geq1-3L\gamma_N=1-\delta_N.                              \tag{39}
\]
Replacing an empty intersection by a full simplex cannot remove the true \(p\), so (39) remains valid for empty arms and zero-probability cells.

By \(\text{thm:sharp-three-facet-endpoints}\), each endpoint is a maximum or minimum over a finite list of affine contrasts
\[
 \ell_r(p)=b_r+\sum_{d\in\mathcal D}\sum_{j=1}^{L-1}g_{rdj}p_{dj}.          \tag{40}
\]
For a lower endpoint, call \(r_*\) pilot-uniformly optimal if \(\ell_{r_*}(\widetilde p)\geq\ell_r(\widetilde p)\) for every \(\widetilde p\in R_N\) and every contrast \(r\) in that endpoint's list.  For an upper endpoint reverse the inequality.  If such an \(r_*\) exists, retain one; otherwise retain the entire finite list.  Thus, throughout \(R_N\), the maximum or minimum over retained contrasts equals the full endpoint map.  This statement remains true at ties: either a retained contrast is optimal throughout the region or the full list is retained.

Condition on the pilot data and on the main-sample arm counts \(m=(m_0,m_1,m_2)\).  Write \(M_{dj}\) for the number of main-sample observations with \((A,Y)=(d,j)\), so that \(m_d=\sum_{j\in\mathcal Y_L}M_{dj}\).  If \(m_d=0\) for an arm on which contrast \(r\) has a nonzero coefficient vector, assign that contrast its full logical range
\[
 \ell_r\!\left(\prod_{d\in\mathcal D}\Delta^{L-1}\right),                 \tag{41}
\]
a compact interval containing \(\ell_r(p)\) surely.  Otherwise define
\[
 \widehat\ell_r=b_r+\sum_{d,j}g_{rdj}\frac{M_{dj}}{m_d},                   \tag{42}
\]
omitting arms whose coefficient vectors vanish.  Let \(K_N\) be the number of retained contrast copies not already handled by (41).  If \(K_N>0\), allocate \(\eta_N=(\alpha-\delta_N)/K_N\) to each; if \(K_N=0\), no testing error is allocated.

For a contrast treated by (42) and \(t\in\ell_r(R_N)\), its null slice
\[
 \mathcal N_r(t)=\{\widetilde p\in R_N:\ell_r(\widetilde p)=t\}            \tag{43}
\]
is nonempty and compact.  Define
\[
 c_r(t)=\min\left\{c\geq0:
 \inf_{\widetilde p\in\mathcal N_r(t)}
 P_{\widetilde p}\left(|\widehat\ell_r-t|\leq c\mid m\right)
 \geq1-\eta_N\right\}.                                                     \tag{44}
\]
This minimum exists.  Conditional on \(m\), the count sample space is finite; for a fixed acceptance set its product-multinomial probability is a polynomial, hence continuous, in \(\widetilde p\), so its infimum is attained on (43).  As \(c\) varies, the acceptance set changes only at the finitely many distances between \(t\) and possible values of (42), and the largest such distance gives acceptance probability one throughout (43).  Thus one of finitely many radii attains the minimum in (44).

Invert these exact tests by
\[
 C_r=\{t\in\ell_r(R_N):|\widehat\ell_r-t|\leq c_r(t)\}.                   \tag{45}
\]
If \(C_r\neq\varnothing\), let \(J_r=\operatorname{cl}\operatorname{conv}(C_r)\); if it is empty, let \(J_r\) be the full logical range (41).  The latter fallback can only enlarge coverage.  Each \(J_r\) is a nonempty compact interval.  On \(\{p\in R_N\}\), setting \(t=\ell_r(p)\) in (43)--(44) yields
\[
 P_p\{\ell_r(p)\in J_r\mid\text{pilot},m\}\geq1-\eta_N.                 \tag{46}
\]
Because all count sample spaces are finite and every operation above is deterministic on them, the resulting interval endpoints are measurable.

Write \(J_r=[L_r,U_r]\).  For a lower endpoint set
\[
 I_d^-=[\max_rL_r,\max_rU_r]\cap[0,1],                                    \tag{47}
\]
and for an upper endpoint set
\[
 I_d^+=[\min_rL_r,\min_rU_r]\cap[0,1].                                    \tag{48}
\]
If either displayed intersection is empty, replace it by \([0,1]\).  Whenever all retained true contrasts lie in their respective \(J_r\), monotonicity of maximum and minimum puts the true endpoints in (47)--(48); since every causal endpoint lies in \([0,1]\), intersection with \([0,1]\) does not remove it.  The Cartesian product of the six intervals is therefore a random closed rectangle \(\mathcal I_{N,1-\alpha}\subseteq[0,1]^6\).

On \(\{p\in R_N\}\), the indicator of failure of any nontrivial retained contrast is bounded by the sum of the individual failure indicators.  Taking conditional expectations and using (46) gives conditional failure probability at most
\[
 K_N\eta_N=\alpha-\delta_N,                                                 \tag{49}
\]
while full-range contrasts have zero failure probability.  Combining (39) and (49), again through the indicator inequality for a union, yields
\[
 P_p\left\{(\underline\theta_d(p),\overline\theta_d(p))_{d\in\mathcal D}
 \notin\mathcal I_{N,1-\alpha}\right\}
 \leq\delta_N+(\alpha-\delta_N)=\alpha.                                   \tag{50}
\]
The construction and bound are conditional on arbitrary realized arm counts and hold for every compatible \(p\), including boundary points.  Taking the infimum over \(P\in\mathcal P_L\) proves the stated finite-sample honesty.  No limit, differentiability, positive-cell, or regular-face condition was used.  In particular, the argument proves no Wald equivalence or scalar efficiency claim.
